\section*{KẾT LUẬN VÀ KIẾN NGHỊ}
\addcontentsline{toc}{section}{KẾT LUẬN VÀ KIẾN NGHỊ}

\subsection*{1. Kết luận}
Nghiên cứu về hoạt động quản lý và quản trị tại Tập đoàn FPT mang lại các kết luận sau:
\begin{enumerate}
    \item FPT đã xây dựng được cơ chế quản lý điều hành linh hoạt và cơ cấu quản trị chuẩn mực, tiệm cận các chuẩn mực quốc tế, tạo nền tảng cho sự tăng trưởng vượt bậc trong hơn 3 thập kỷ.
    \item Tuy nhiên, quy mô nhân sự ngày càng lớn đặt ra bài toán về độ trễ điều hành, sự chồng chéo thẩm quyền giữa thế hệ sáng lập và lãnh đạo kế cận, đòi hỏi phải tiếp tục tinh gọn và phẳng hóa tổ chức.
    \item Cuộc cách mạng Trí tuệ Nhân tạo đang mở ra cơ hội đột phá cho FPT trong việc tối ưu hóa năng suất lao động và tự động hóa quy trình quản lý, nhưng cũng đặt ra những thách thức lớn về an toàn dữ liệu, tính minh bạch và đạo đức công nghệ trong quản trị.
\end{enumerate}

\subsection*{2. Kiến nghị đối với Tập đoàn FPT}
\begin{itemize}
    \item \textbf{Về quản lý:} Đẩy mạnh trao quyền và phân cấp trách nhiệm thực thi; xóa bỏ rào cản chia sẻ dữ liệu giữa các khối ngành để khai thác tối đa sức mạnh hệ sinh thái.
    \item \textbf{Về quản trị:} Nâng cao tính độc lập thực chất của thành viên HĐQT độc lập; hoàn thiện khung chính sách quản trị AI (AI Governance Framework) nhằm giảm thiểu rủi ro pháp lý và bảo mật.
    \item \textbf{Về phát triển nguồn nhân lực:} Xây dựng chiến lược đào tạo AI phổ cập cho toàn bộ đội ngũ quản lý các cấp để thay đổi tư duy điều hành theo thời gian thực.
\end{itemize}
